\clearpage
\section{The literal copying rule and preservation of shared stores}
\label{sec:copying-rule}
Figure~17 (PDF p.~37) now has a concrete finite-graph implementation, for
an arbitrary term contraction tree $C$ with $n+1$ inputs and an arbitrary
copyable hypernet $H:\star\Rightarrow\star^k\otimes\diamond^h$.
The selected occurrence is input zero; the other $n$ inputs keep their order.
The rewrite duplicates $H$, retains $C$ behind the second copy, weakens the
consumed input of $C$, and inserts the paper's two output distributors.
Neither the proposition's terms nor the translation clauses are changed.

This milestone proves typing and finiteness by constructing actual nets,
recursive safety, disjointness of copied edge identities, preservation of
literal box contents, carrier-renaming laws and forward padding laws. Local
copy-and-return executions and shared-reference resolution are checked.
It does not yet prove independence from different contraction-tree matches
in one ambient graph, arbitrary focused-context matching, or determinism and
refocusing of the complete copying UAM. Those remain in Appendix~B.1.

\Needspace{12\baselineskip}
\subsection{Retaining the matched contraction tree}

\termdef{copying-residual-tree}
The residual tree is the original matched tree with its first input closed
by a term weakening. Its interface is $\star^n\Rightarrow\star$.
The construction uses a two-edge frame: a weakening and a tree hole of
signature $\star^{n+1}\Rightarrow\star$. One internal term vertex connects
the weakening output to input zero of the hole. The remaining hole inputs
are the external inputs in increasing order; its output is the external
output. The label and incidence functions specify this wiring directly.

\lean{Copying/CopyFrames}{9}{9}{CopyRule.treeHole}

\lean{Copying/CopyFrames}{10}{12}{CopyRule.residualLabel}

\lean{Copying/CopyFrames}{13}{15}{CopyRule.residualSource}

\lean{Copying/CopyFrames}{16}{18}{CopyRule.residualTarget}

The frame is a finite typed net with a Boolean edge carrier and one internal
vertex. Its constructor checks each source and target for duplicate freedom,
unique incidence, complete coverage and agreement with the primitive signature.
The two edge cases discharge uniqueness; coverage splits vertices into the
single internal vertex, external inputs and external output. The ordered
finite input range gives both the typing equation and input injectivity.

\lean{Copying/CopyFrames}{21}{55}{CopyRule.residualFrame}

Filling the tree slot with $C$ and the other slot with the weakening corolla
produces the residual graph. Relabelling primitive graphs into allowance $d$
produces both the original tree and the residual tree at that allowance;
it does not alter any graph incidence. The two edge injections distinguish
the new weakening from every original tree edge. Their label theorem follows
by unfolding expansion. Increasing the allowance commutes literally with
both primitive-only constructions, by the two depth cases.

\lean{Copying/CopyFrames}{57}{62}{CopyRule.residualFill}

\lean{Copying/CopyFrames}{64}{66}{CopyRule.tree}

\lean{Copying/CopyFrames}{68}{71}{CopyRule.residual}

\lean{Copying/CopyFrames}{73}{73}{CopyRule.residualWeakening}

\lean{Copying/CopyFrames}{74}{74}{CopyRule.residualTreeEdge}

\lean{Copying/CopyFrames}{76}{78}{CopyRule.residual_labels}

\lean{Copying/CopyFrames}{80}{80}{CopyRule.tree_raise}

\lean{Copying/CopyFrames}{81}{81}{CopyRule.residual_raise}

The tree-vertex map is the expansion insertion map. The selected input is
exactly the weakening's target, not an erased wire. Original input $i+1$
becomes external input $i$; this follows by indexing the explicit finite
input list. Every original source and target list is mapped entrywise by
the same insertion map. Thus the entire ordered incidence of $C$ survives,
not merely its number of leaves or its output type.

\lean{Copying/CopyExamples}{8}{9}{CopyRule.residualTreeVertex}

\lean{Copying/CopyExamples}{12}{14}{CopyRule.selected_input_weakened}

\lean{Copying/CopyExamples}{17}{21}{CopyRule.residual_input_order}

\lean{Copying/CopyExamples}{24}{25}{CopyRule.residual_sources}

\lean{Copying/CopyExamples}{27}{28}{CopyRule.residual_targets}

\Needspace{12\baselineskip}
\subsection{The two rule graphs and their side condition}

\termdef{copying-intrinsic-rule}
Write $R(C)$ for the residual tree. Before copying the graph is $C;H$.
After copying it is
\[
  \bigl(H\otimes(R(C);H)\bigr);\chi;
  \bigl(D^{\star}_{k,2}\otimes D^{\diamond}_{h,2}\bigr),
\]
where $\chi$ is the middle-block exchange of Figure~17. Each copy of $H$
has term outputs followed by atom outputs; the exchange groups the two term
blocks before the two atom blocks. It uses the same ordered output bijection
as the original binary translation. The external input order is the active
input followed by the $n$ residual inputs. Interface equality transport
changes only length indices and preserves edge and internal carriers.
The full canonical distributors, including their weakening edges, remain.

\lean{Copying/CopyRule}{10}{19}{CopyRule.pairBoundary}

\lean{Copying/CopyRule}{22}{27}{CopyRule.copies}

\lean{Copying/CopyRule}{29}{31}{CopyRule.before}

\lean{Copying/CopyRule}{33}{36}{CopyRule.after}

The intrinsic rule relates these two concrete graphs only when $H$ is
copyable in Definition~10.2's existing formalisation. This includes instances,
constants, active operations and lambda payloads. It does not replace
copyability by stability. The rule-characterisation theorem says precisely
that any rule witness supplies such a tree and copyable payload, with the
two graph equalities. Its proof eliminates or constructs the sole rule
constructor. The transition is separate from the linear compute relation.

\lean{Copying/CopyRule}{39}{44}{CopyRule.Rule}

\lean{Copying/CopyRule}{46}{50}{CopyRule.rule_iff}

The left and right edge injections use different Boolean component tags.
Their labels are exactly the original labels, proved by unfolding the
constructors at zero and successor allowance. The injections have disjoint
images: equality would force the distinct Boolean tags to be equal.
In particular a box label, including its complete body, is preserved exactly.

\lean{Copying/CopyRule}{52}{54}{CopyRule.leftEdge}

\lean{Copying/CopyRule}{55}{57}{CopyRule.rightEdge}

\lean{Copying/CopyRule}{60}{63}{CopyRule.copy_labels}

\lean{Copying/CopyRule}{66}{73}{CopyRule.copies_disjoint}

\Needspace{12\baselineskip}
\subsection{Safety throughout nested boxes}

Let $P$ be a primitive-label predicate. If $P$ holds on term weakening and
contraction, it holds throughout the tree and its residual. For the original
tree this is the tree's label classification; for the residual, split an
edge into the inserted weakening or an original tree edge. If $H$ also
satisfies $P$ recursively, serial assembly proves safety of the redex.
For the reduct, require $P$ on weakening and contraction of both sorts.
Split parallel components, use residual safety for the second component,
then use the distributor and serial-assembly label laws. No assumption that
box bodies are passive is introduced.

\lean{Copying/CopySafety}{8}{11}{CopyRule.tree_labels}

\lean{Copying/CopySafety}{13}{17}{CopyRule.residual_allLabels}

\lean{Copying/CopySafety}{19}{22}{CopyRule.before_allLabels}

\lean{Copying/CopySafety}{24}{33}{CopyRule.after_allLabels}

Taking $P$ to exclude focus labels proves focus freedom of the reduct;
taking it to exclude hole labels proves hole freedom. These are recursive
properties, so they also exclude forbidden labels inside copied boxes.
The frame's temporary tree hole has been filled and does not remain.

\lean{Copying/CopySafety}{35}{37}{CopyRule.focusFree_after}

\lean{Copying/CopySafety}{39}{41}{CopyRule.holeFree_after}

\Needspace{12\baselineskip}
\subsection{Local copy execution}

\termdef{copying-copy-subsystem}
A local state consists of a concrete graph at the common rule interface and
a focus mode with a vertex. Its transition relation contains intrinsic
search and the copy transition from rewrite focus at input zero to search
focus at that same external input on the reduct. Finality means backtracking
at an external input. A final state has no transition: search is impossible
by search finality, while copying requires a different mode. This defines a
search/copy subsystem with a proved final-normal law. It is not a claim that
the entire copying UAM, including all compute rules and contexts, is finished.

\lean{Copying/CopyExecution}{8}{9}{CopyRule.State}

\lean{Copying/CopyExecution}{11}{15}{CopyRule.Step}

\lean{Copying/CopyExecution}{17}{25}{CopyRule.final_normal}

\lean{Copying/CopyExecution}{27}{30}{CopyRule.machine}

For an instance payload, the active copy has an instance edge at input zero,
so its next search transition backtracks there. For any constant payload,
the same argument uses the zero-argument passive-operation search rule.
The incidence calculation unfolds the parallel boundary and serial insertion;
the active input remains exactly position zero. Composing one intrinsic
copy transition and this return proves a two-step local execution in each
case. These runs begin at rewrite focus, rather than claiming an unproved
whole-program path from the original source term.

\lean{Copying/CopyExecution}{33}{40}{CopyRule.instance_returns}

\lean{Copying/CopyExecution}{43}{51}{CopyRule.constant_returns}

\lean{Copying/CopyExecution}{54}{58}{CopyRule.instance_run}

\lean{Copying/CopyExecution}{60}{64}{CopyRule.constant_run}

\Needspace{12\baselineskip}
\subsection{References share their original cell}

\termdef{copying-shared-store}
To test reference copying, attach a single original store to the atom output
of the copy reduct for an instance payload. This test is at allowance zero
and allows any residual term tree. The two designated instance edges are
nested left/right injections; the designated store is the right serial
component. The atom distributor has the actual canonical two-use tree:
a weakening, a lower contraction and its root contraction. Its edge selector
chooses these three occurrences. Both instance labels remain instance labels,
and the store label remains a store label.

\lean{Copying/CopyStore}{9}{11}{CopyRule.sharedStore}

\lean{Copying/CopyStore}{13}{14}{CopyRule.sharedLeft}

\lean{Copying/CopyStore}{15}{16}{CopyRule.sharedRight}

\lean{Copying/CopyStore}{17}{18}{CopyRule.sharedCell}

\lean{Copying/CopyStore}{19}{20}{CopyRule.sharedContraction}

\lean{Copying/CopyStore}{22}{24}{CopyRule.shared_instance_labels}

\lean{Copying/CopyStore}{26}{27}{CopyRule.shared_cell_label}

Both instance outputs enter the same lower atom contraction, at its first
and second input respectively. That contraction feeds the root contraction,
which feeds the single original store. Each hop proof supplies the actual
shared target/source vertex and checks its ordered incidence membership.
Composing these hops gives exactly three atom-routing steps from either
copy to the same store. Store-trace uniqueness then rules out any different
successful destination or hop count. The store is shared; copying has not
allocated a new cell for either occurrence.

\lean{Copying/CopyStore}{30}{36}{CopyRule.shared_left_hop}

\lean{Copying/CopyStore}{38}{44}{CopyRule.shared_right_hop}

\lean{Copying/CopyStore}{46}{52}{CopyRule.shared_root_hop}

\lean{Copying/CopyStore}{54}{60}{CopyRule.shared_store_hop}

\lean{Copying/CopyStore}{63}{68}{CopyRule.shared_resolves}

\lean{Copying/CopyStore}{71}{76}{CopyRule.shared_resolves_unique}

\Needspace{12\baselineskip}
\subsection{Renaming and allowance transport}

Boundary equality reindexing preserves a deep isomorphism, by eliminating
the two interface equalities. It commutes with raising by the same argument.
An ordered isomorphism between two matched tree carriers lifts to the
copying hierarchy. For residual trees, use componentwise expansion
isomorphism on the two-edge frame, with identity on the weakening.
Parallel component isomorphisms then transport the two copies, and serial
composition transports both complete rule sides. The payload is allowed a
deep isomorphism, which also renames recursively inside boxes.

\lean{Copying/CopyTransport}{9}{11}{CopyRule.reboundaryIso}

\lean{Copying/CopyTransport}{13}{15}{CopyRule.raise_reboundary}

\lean{Copying/CopyTransport}{17}{18}{CopyRule.tree_congr}

\lean{Copying/CopyTransport}{20}{25}{CopyRule.residual_congr}

\lean{Copying/CopyTransport}{27}{33}{CopyRule.copies_congr}

\lean{Copying/CopyTransport}{35}{37}{CopyRule.before_congr}

\lean{Copying/CopyTransport}{39}{41}{CopyRule.after_congr}

Raising commutes with the parallel constructor and boundary reindexing.
The residual serial component commutes up to the already proved serial
padding isomorphism; the residual primitive graph itself raises literally.
These facts prove padding for the paired copies. Serial padding gives the
redex law, and the paired-copy law followed by distributor padding gives the
reduct law. These are explicit operational construction proofs, not an
inference from determinism.

\lean{Copying/CopyTransport}{43}{52}{CopyRule.copies_raise}

\lean{Copying/CopyTransport}{54}{57}{CopyRule.before_raise}

\lean{Copying/CopyTransport}{59}{64}{CopyRule.after_raise}

The representation closure of a copy pair consists of a tree and a copyable
payload, together with deep isomorphisms from the supplied two graphs to the
literal redex and reduct. Every raw rule supplies such a witness. Independent
renamings of its two sides preserve this closure by inverse composition.
Raising transports the two witnesses using the preceding padding laws and
the proved preservation of copyability. Induction gives every finite padding
amount. This is forward preservation; no new converse execution theorem is
being claimed for the combined UAM.

\lean{Copying/CopyTransport}{67}{68}{CopyRule.UpToIso}

\lean{Copying/CopyTransport}{70}{73}{CopyRule.rule_upToIso}

\lean{Copying/CopyTransport}{75}{78}{CopyRule.upToIso_congr}

\lean{Copying/CopyTransport}{80}{84}{CopyRule.upToIso_raise}

\lean{Copying/CopyTransport}{86}{90}{CopyRule.upToIso_pad}

\Needspace{12\baselineskip}
\subsection{Operational regressions and remaining obligations}

The lambda regression applies the rule to an arbitrary original lambda
payload, with no stability restriction on its body. For every box-labelled
edge, both copied labels equal the original box label by transitivity with
the general label theorem. This includes nested boxes without translating
or normalising their contents. A separate application regression checks
that active application heads also satisfy copyability.

\lean{Copying/CopyExamples}{31}{34}{CopyRule.lambda_rule}

\lean{Copying/CopyExamples}{37}{43}{CopyRule.lambda_box_preserved}

\lean{Copying/CopyExamples}{46}{48}{CopyRule.application_rule}

The weakening-free binary tree supplies a noncanonical matched tree.
Its residual contains the new weakening even though the original tree had
none. Copying a reference with this match gives two distinct instance edges
that resolve to the same cell in three hops, by disjointness and the shared
resolution theorem. These checks do not yet constitute repeated closure
application or execution of a divergent callback in the complete language;
those require the missing application rule and machine integration.

\lean{Copying/CopyExamples}{52}{54}{CopyRule.binary_residual_weakening}

\lean{Copying/CopyExamples}{57}{66}{CopyRule.binary_shared_reference}

The build checks 318 jobs. The axiom audits check 6 abstract, 477 graph and
528 copying results and use only the standard Lean logical axioms.
Estimated completion of Proposition~12.2 is about 63\% by mathematical effort,
with broad uncertainty of 48--73\% (26 October 2026). The copying construction
is now proved, but matching independence, the remaining compute rules,
full-machine integration and the eleven template robustness proofs remain
substantial. Appendix~B.1 removes the completed local copy construction
from its pending list while retaining these exact obligations.

